\documentclass[12pt, letterpaper]{article}
\usepackage{graphicx}
\usepackage[margin=0.5in]{geometry} 
\usepackage{array}
\usepackage{amsmath}
\renewcommand{\familydefault}{\sfdefault}



\begin{document}

\begin{center}
    {\Large\bfseries Assignment 1: Order of Operations}\\[4pt]
    {\large\bfseries DO NOT WRITE ON THIS PAPER}\\[4pt]
\end{center}



\fbox{%
  \begin{minipage}[t][0.40\textheight][t]{.95\textwidth}
    \begin{center}
        \textbf{Part A}
    \end{center}

    Roll a die three times and use the rolls to replace A, B, and C in the following expression.  

    \[(A + 3)^B + 4 \times C - 10 + 5 - 35 \div 4^0\]


    \emph{Example: If you rolled a 4, 5, and 1, your expression would be: \((4 + 3)^5 + 4 \times 3 - 10 + 5 - 35 \div 4^0\)} 

    \vspace{1.5em}


    Then, simplify the expression using the order of operations.  You should have four clear sections of steps each corresponding to a letter in PEMDAS, organized into a table like the one shown below.   


    



    \begin{center}
        \begin{tabular}{|c|c|c|}
        \hline
        \textbf{Letter} & \textbf{Name of Operation} & \textbf{Calculations} \\
        \hline
        P &  &  \\
        \hline
        E &  &  \\
        \hline
        \makebox[2em][l]{M\hspace{1em}\raisebox{-1em}{D}} &  &  \\
        \hline
        \makebox[2em][l]{A\hspace{1em}\raisebox{-1em}{S}} &  &  \\
        \hline
        \end{tabular}
    \end{center}

    \vfill
    \hfill
  \end{minipage}%
}%


\fbox{%
  \begin{minipage}[t][0.46\textheight][t]{.95\textwidth}
    \begin{center}
        \textbf{Part B}
    \end{center}

    Griffin was tired and wrote the following expression and simplified it as follows:  

    \vspace{0.5em}
    \[
    \begin{aligned}
    \text{Given Problem:} \quad & (3-2)^2 - (15 \div 3) \\
    \text{Step 1:} \quad & = (3-2)^2 - 5 \\
    \text{Step 2:} \quad & = 3^2 - 2^2 - 5 \\
    \text{Step 3:} \quad & = 9 - 4 - 5 \\
    \text{Step 4:} \quad & = 5 - 5 \\
    \text{Answer:} \quad & = 0
    \end{aligned}
    \]

    \vspace{0.5em}
    What mistake did Griffin make and at what step did it happen?

    Correct any mistakes Griffin made and give the true value of the expression. 

    \vfill
    \hfill
  \end{minipage}%
}%

\newpage

\fbox{%
  \begin{minipage}[t][0.40\textheight][t]{.95\textwidth}
    \begin{center}
        \textbf{Part C}
    \end{center}

    Variables are reprsentations of numbers: they are place holders. Knowing this and order of operations are the following expressions equivalent? Explain your reasoning.    

    \vspace{2em}

    \[
    \text{Is }(x + 4)^2  \text{  equivalent to  } x^2 + 16 \text{?}
    \]

    \vspace{2em}
    \emph{Hint: You can plug in any number for x and simplify. Try at least two different numbers to see if the expressions are really the same.}

    \vfill
    \hfill
  \end{minipage}%
}

%\setlength{\fboxsep}{0pt} 
\fbox{%
  \begin{minipage}[t][0.46\textheight][t]{.95\textwidth}
    \begin{center}
        \textbf{Question 1} 
    \end{center}



    \textbf{Big Question:} What is the Order of Operations? Why is it important that we follow them?
    \vspace{2em}

    Use the example below to help support your answer as to why they are important.

    \[
    5 + 9 \div 3
    \]


    \vfill
    \hfill
  \end{minipage}%
}%

\newpage

\fbox{%
  \begin{minipage}[t][0.42\textheight][t]{.7\textwidth}
    \begin{center}
        \textbf{Question 2 Work}
    \end{center}
    \vspace{1em}

    Evaluate the following expression.

        \[
        2 \times [1 + (12 + 3) \div 3]
        \]

   

    \hfill
    \vfill
  \end{minipage}%
}%
\fbox{%
  \begin{minipage}[t][0.42\textheight][t]{0.24\textwidth}
    \begin{center}
      {\bfseries Question 2 Answer}
      
        \vspace*{\fill}
      
      {\Large Write your answer in your Class Notes}
    \end{center}
    \vspace*{\fill}  % pushes content up, centers vertically
  \end{minipage}%
}

\fbox{%
  \begin{minipage}[t][0.42\textheight][t]{.7\textwidth}
    \begin{center}
        \textbf{Question 3 Work}
    \end{center}

    Joe is buying ribbon for wrapping presents. A large present requires 8.25 feet of ribbon, a medium present requires 5.5 feet, and a small present requires 3.25 feet.
    \vspace{1em}
    
    If Joe is wrapping 9 large presents, 10 medium presents, and one small present, how much ribbon will he need?

    \vspace{1em}

    Write your answer using a complete sentence. 



  \end{minipage}%
}%
\fbox{%
  \begin{minipage}[t][0.42\textheight][t]{0.24\textwidth}
  % pushes content down
    \begin{center}
      {\bfseries Question 3 Answer}
      
        \vspace*{\fill}
      
      {\Large Write your answer in your Class Notes}
    \end{center}
    \vspace*{\fill}  % pushes content up, centers vertically
  \end{minipage}%
}

\newpage
\begin{center}
    \textbf{Additional Space}
\end{center}

\vspace{0.4em}

\noindent\makebox[\textwidth][c]{%
\fbox{%
  \rule{0pt}{0.20\textheight}\rule{7.35in}{0pt}%
}%
}

\vspace{0.45em}

\noindent\makebox[\textwidth][c]{%
\fbox{%
  \rule{0pt}{0.20\textheight}\rule{7.35in}{0pt}%
}%
}

\vspace{0.45em}

\noindent\makebox[\textwidth][c]{%
\fbox{%
  \rule{0pt}{0.20\textheight}\rule{7.35in}{0pt}%
}%
}

\vspace{0.45em}

\noindent\makebox[\textwidth][c]{%
\fbox{%
  \rule{0pt}{0.20\textheight}\rule{7.35in}{0pt}%
}%
}

\end{document}
